% BCT Zenodo Companion Letter — Appendix JH: The Hopfion Interior
% M. R. Cabrié — March 2026
% Format: revtex4-2, aps/prl, twocolumn — identical to BCT Letter series

\documentclass[aps,prl,twocolumn,superscriptaddress,floatfix]{revtex4-2}

% ── Packages ────────────────────────────────────────────────────
\usepackage{amsmath,amssymb,amsfonts}
\usepackage{graphicx}
\usepackage{dcolumn}
\usepackage{bm}
\usepackage{hyperref}
\usepackage{booktabs}
\usepackage{xcolor}

% ── Hyperref ────────────────────────────────────────────────────
\hypersetup{
  colorlinks=true,
  linkcolor=blue!70!black,
  citecolor=blue!70!black,
  urlcolor=blue!70!black
}

% ── BCT standard macros ─────────────────────────────────────────
\newcommand{\roct}{r_{\mathrm{oct}}}
\newcommand{\rtet}{r_{\mathrm{tet}}}
\newcommand{\az}{\alpha_0}
\newcommand{\mP}{m_{\mathrm{P}}}
\newcommand{\lP}{\ell_{\mathrm{P}}}
\newcommand{\LQCD}{\Lambda_{\mathrm{QCD}}}
\newcommand{\BCT}{\textsc{bct}}
\newcommand{\OHC}{\textsc{ohc}}
\newcommand{\Psiint}{\Psi_2}
\newcommand{\Psiext}{\Psi_1}
\newcommand{\kz}{\kappa_0}
\newcommand{\HopfN}{H}
\newcommand{\vph}{v_{\mathrm{ph}}}
\newcommand{\Rsph}{R_{\mathrm{s}}}

% ════════════════════════════════════════════════════════════════
\begin{document}

\title{%
  BCT Letter~36: The Octet-Hopfion Condensate\\
  Topology of the \BCT\ Interior Condensate\\
  \normalsize{Companion to Appendix~JH: The Hopfion Interior}
}

\author{M.~R.~Cabri\'{e}}
\email{ZeroFreeParameters@gmail.com}
\affiliation{Independent Researcher; ORCID: 0009-0007-9561-9859}

\date{Dated: March 2026 \quad BCT Letter~36 \quad [Patent Pending]}

% ── Abstract ────────────────────────────────────────────────────
\begin{abstract}
This letter accompanies the Zenodo upload of BCT Appendix~JH:
\textit{The Hopfion Interior} (BCT Programme, March 2026).
It states the principal results, identifies the central new concept
introduced (the Octet-Hopfion Condensate, \OHC), locates the appendix
within the ongoing \BCT\ publication programme, and provides a complete
citation block.
The \OHC\ is the ground-state field configuration of the \BCT\ interior
condensate $\Psiint$, characterised by Hopf invariant $\HopfN=1$,
and representing the unique smooth solution to the Gross--Pitaevskii
equation with $\az$~Josephson boundary conditions inside each \BCT\
Planck-scale sphere.
Six principal results are established: topological quantisation, spinor
topology, the \BCT--twistor correspondence, three fermion generations
from $D_4$ triality, superluminal interior phase velocity, and the
electron-mass mechanism.
\end{abstract}

\maketitle

% ════════════════════════════════════════════════════════════════
\section{Purpose of This Letter}
\label{sec:purpose}
% ════════════════════════════════════════════════════════════════

This companion letter is provided with BCT Appendix~JH uploaded to
Zenodo on 12~March 2026.
It serves three purposes:
\begin{enumerate}
\item[(a)] To state the principal results of Appendix~JH in condensed
  form, suitable for readers deciding whether to read the full appendix.
\item[(b)] To introduce and formally define the \textbf{Octet-Hopfion
  Condensate (\OHC)}, which is the central new concept of Appendix~JH.
\item[(c)] To locate Appendix~JH within the ongoing \BCT\ publication
  programme and provide a complete citation block for the published
  open-access record.
\end{enumerate}

% ════════════════════════════════════════════════════════════════
\section{The Octet-Hopfion Condensate: Definition}
\label{sec:ohc}
% ════════════════════════════════════════════════════════════════

The \BCT\ Superfluid Lattice Model~\cite{Monograph} proposes that the
quantum vacuum is a Body-Centred Tetragonal superfluid lattice of
Planck-scale spherical oscillators with axial ratio $c/a=\sqrt{2}$.
Each sphere of radius $\Rsph = \tfrac{1}{2}\lP$ contains an interior
condensate $\Psiint$ coupled to the exterior void condensate $\Psiext$
through a Josephson junction of strength
\begin{equation}
  \az = \frac{\roct\cdot\rtet}{\pi} = 0.00740806\ldots,
  \label{eq:alpha0}
\end{equation}
where
\begin{align}
  \roct &= \tfrac{\sqrt{2}-1}{2} = 0.207107\ldots, \nonumber\\
  \rtet &= \tfrac{\sqrt{6}-2}{4} = 0.112372\ldots
  \label{eq:radii}
\end{align}
are the octahedral and tetrahedral void radii in units of the
nearest-neighbour distance.

Appendix~JH establishes that the natural ground-state topology of
$\Psiint$ is the Hopf fibration $\eta: S^3\to S^2$, with Hopf invariant
$\HopfN=1$ (the generator of $\pi_3(S^2)=\mathbb{Z}$).
We define:

\begin{quote}
\textbf{Octet-Hopfion Condensate (\OHC):}
The ground-state field configuration of the \BCT\ interior condensate
$\Psiint$, characterised by Hopf invariant $\HopfN=1\in\mathbb{Z}$.
The name reflects the octahedral-tetrahedral \BCT\ void geometry
\emph{(Octet)} and the Hopf-fibration internal topology
\emph{(Hopfion)}.
The \OHC\ is forced by the $\az$~Josephson boundary condition, is
topologically protected, and is the unique smooth ground state consistent
with the $D_4$~symmetry of the \BCT\ lattice.
\end{quote}

This definition is introduced for the first time in this appendix.
The formal abbreviation \OHC\ is used in all subsequent \BCT\
publications.

% ════════════════════════════════════════════════════════════════
\section{Principal Results}
\label{sec:results}
% ════════════════════════════════════════════════════════════════

Appendix~JH establishes six principal results, summarised here.

\subsection{Result 1: OHC Ground State}

The interior condensate $\Psiint$ ground state is the \OHC\ with
$\HopfN=1$, forced by the $\az=0.00740806$ Josephson coupling.
$\HopfN$ is conserved and cannot change under any smooth deformation
of finite energy.

\subsection{Result 2: Quantisation from Topology}

The Hopf invariant $\HopfN\in\mathbb{Z}$ is integer-valued by algebraic
topology (Bott--Milnor theorem~\cite{Bott1958}).
The energy spectrum of the interior condensate is
\begin{align}
  E_n &= n\,\varepsilon_0, \quad n\in\mathbb{Z}, \nonumber\\
  \varepsilon_0 &= \frac{\hbar^2\kz^2}{2m\Rsph^2},
  \label{eq:spectrum}
\end{align}
with $\kz=3.39$ from Appendix~J~\cite{AppJ}.
The canonical commutation relation $[\hat{x},\hat{p}]=i\hbar$ is a
consequence of the \OHC\ topology, not a postulate.

\subsection{Result 3: Half-Integer Spin}

The Hopf fibration $S^3\to S^2$ realises the double cover
$\mathrm{SU}(2)\to\mathrm{SO}(3)$.
Fermions threading the \OHC\ interior accumulate a phase of $\pi$ per
traversal ($4\pi$ periodicity).
Half-integer spin is a topological consequence of the \OHC\ structure,
not an input to \BCT.

\subsection{Result 4: BCT--Twistor Correspondence}

Each \BCT\ Planck sphere with \OHC\ Hopf charge $\HopfN=n$ corresponds
to a Penrose twistor~\cite{Penrose1967,PR1986} of helicity $h=n/2$.
The interior Hopf fibration \emph{is} the twistor fibration.
The \BCT\ lattice of \OHC\ spheres is Penrose twistor space at the
Planck scale.

\subsection{Result 5: Three Generations from $D_4$ Triality}

The $D_4$~triality outer automorphism $S_3$ of the \BCT\ root lattice
cyclically permutes the three $\mathrm{Spin}(8)$ representations
$8_v, 8_s, 8_c$.
The three Standard Model fermion generations correspond to the three
triality-related \OHC\ states.
The \BCT\ framework is formally an octonionic theory.

\subsection{Result 6: Superluminal Interior Phase Velocity}

The \OHC\ ground-state phase velocity is
\begin{equation}
  \vph = \frac{\kz}{\Rsph} = \frac{3.39}{0.5}\,c \approx 6.78\,c.
  \label{eq:vph}
\end{equation}
This is a phase velocity, causally screened behind the $\az$~Josephson
shell.
The velocity mismatch $\vph/c\approx 6.78$ between the
subluminal exterior ($c$) and the superluminal \OHC\ interior is
the structural mechanism behind the Kondo-like electron-mass
suppression~\cite{Letter19,Letter34}.

% ════════════════════════════════════════════════════════════════
\section{Publication Context}
\label{sec:context}
% ════════════════════════════════════════════════════════════════

BCT Appendix~JH is Record~27 of the \BCT\ publication programme,
which targets 50~open-access Zenodo records in total.
At the time of this upload, 26~records have been published
(Records~1--26, covering the BCT Monograph, Appendices Vol.~1--2,
and Letters~1--22 plus selected letters).
The remaining 24~records (Letters~23--35 and additional appendices)
will be released over the coming days.

The complete \BCT\ programme comprises 110~sub-1\% Standard Model
predictions, 40~sub-0.1\% predictions, zero free parameters, and three
geometric inputs ($\roct$, $\rtet$, $\LQCD=220\,\mathrm{MeV}$).

\subsection*{Previously Published Records Relevant to This Appendix}

\begin{itemize}
\item \textbf{BCT Monograph} (Record~2):
  doi:\href{https://doi.org/10.5281/zenodo.18884976}{10.5281/zenodo.18884976}

\item \textbf{BCT Appendices Vol.~1 \& 2} (Records~3--4):
  doi:\href{https://doi.org/10.5281/zenodo.18885133}{10.5281/zenodo.18885133},
  \href{https://doi.org/10.5281/zenodo.18885472}{10.5281/zenodo.18885472}

\item \textbf{BCT Letter~18} --- $D_4$ Uniqueness Theorem (Record~14):
  doi:\href{https://doi.org/10.5281/zenodo.18957202}{10.5281/zenodo.18957202}

\item \textbf{BCT Letter~19} --- Electron Mass (Record~7):
  doi:\href{https://doi.org/10.5281/zenodo.18905765}{10.5281/zenodo.18905765}

\item \textbf{BCT Letter~34} --- All Three SM Scales (Record~13):
  doi:\href{https://doi.org/10.5281/zenodo.18908349}{10.5281/zenodo.18908349}
\end{itemize}

% ════════════════════════════════════════════════════════════════
\section{Citation Block for This Record}
\label{sec:citation}
% ════════════════════════════════════════════════════════════════

\begin{quote}
M.~R.~Cabri\'{e},
``BCT Letter~36: The Octet-Hopfion Condensate ---
Companion Letter to BCT Appendix~JH,''\
BCT Programme Record~28 (2026),
Zenodo, \href{https://doi.org/10.5281/zenodo.18974754}{doi:10.5281/zenodo.18974754}.
\end{quote}

\noindent Companion appendix (Record~27):

\begin{quote}
M.~R.~Cabri\'{e},
``BCT Appendix~JH: The Hopfion Interior ---
Topology of the Interior Condensate and the Octet-Hopfion Condensate,''
BCT Programme Record~27 (2026),
Zenodo, \href{https://doi.org/10.5281/zenodo.18975018}{doi:10.5281/zenodo.18975018}.
\end{quote}

For the full programme, see ORCID:~0009-0007-9561-9859 and
\href{https://zenodo.org/search?q=cabri\%C3\%A9}{zenodo.org/search?q=cabri\'{e}}.

% ════════════════════════════════════════════════════════════════
\begin{thebibliography}{99}

\bibitem{Monograph}
M.~R.~Cabri\'{e},
``The BCT Superfluid Lattice Model,''
BCT Monograph (2026),
\href{https://doi.org/10.5281/zenodo.18884976}{doi:10.5281/zenodo.18884976}.

\bibitem{AppJ}
M.~R.~Cabri\'{e},
``BCT Appendix J: The Sphere Interior Condensate Action $S[\Psi_2]$,''
in \textit{BCT Appendices Vol.~2} (2026),
\href{https://doi.org/10.5281/zenodo.18885472}{doi:10.5281/zenodo.18885472}.

\bibitem{Letter18}
M.~R.~Cabri\'{e},
``BCT Letter~18: The $D_4$ Uniqueness Theorem,''
BCT Programme (2026),
\href{https://doi.org/10.5281/zenodo.18957202}{doi:10.5281/zenodo.18957202}.

\bibitem{Letter19}
M.~R.~Cabri\'{e},
``BCT Letter~19: The Electron Mass from \BCT\ Vacuum Geometry,''
BCT Programme (2026),
\href{https://doi.org/10.5281/zenodo.18905765}{doi:10.5281/zenodo.18905765}.

\bibitem{Letter34}
M.~R.~Cabri\'{e},
``BCT Letter~34: All Three Fundamental SM Scales from \BCT\ Geometry,''
BCT Programme (2026),
\href{https://doi.org/10.5281/zenodo.18908349}{doi:10.5281/zenodo.18908349}.

\bibitem{Penrose1967}
R.~Penrose,
``Twistor Algebra,''
\textit{J. Math. Phys.} \textbf{8}, 345 (1967).

\bibitem{PR1986}
R.~Penrose and W.~Rindler,
\textit{Spinors and Space-Time, Vol.~2},
Cambridge University Press (1986).

\bibitem{Bott1958}
R.~Bott and J.~Milnor,
``On the Parallelizability of the Spheres,''
\textit{Bull. Amer. Math. Soc.} \textbf{64}, 87 (1958).

\bibitem{Faddeev1997}
L.~D.~Faddeev and A.~J.~Niemi,
``Stable Knot-Like Structures in Classical Field Theory,''
\textit{Nature} \textbf{387}, 58 (1997).

\end{thebibliography}

\end{document}
